\subsection{Alternative Braid-to-Manifold Identifications}
\label{subsec:alternative-identification}

We retain the notation and conventions of Section~\ref{sec:mtc} (modular tensor category \(\mathcal{C}_p = \mathrm{SU}(2)_{p-2}\)) and Section~\ref{sec:braid} (braid encoding of periodic orbits on \(T_p\)). The remark in Subsection~\ref{subsec:verlinde-reduction} identifies two independent obstructions to a rigorous proof of \(\mathrm{RT}_{\mathcal{C}_p}(\beta) = p^{-n}\): (i) the conjectural collapse of the Verlinde state sum to the single Gauss-sum term \(G_p\), and (ii) the failure of the identity \(G_p = \tau_p\) (with \(\tau_p = e^{\pi i/4} p^{-1/2}\)) under the standard MTC data. The present subsection explores the first of the two proposed remedies: a \emph{fundamentally different braid-to-manifold identification} that alters the Verlinde matrix element appearing in the reduced expression.

\subsubsection{What the alternative construction must achieve}
A braid-to-manifold identification consists of three data:
\begin{enumerate}
\item An explicit word \(\beta \in B_n\) (or, more generally, a framed link diagram \(\mathcal{D}_\beta\)) canonically associated to each length-\(n\) periodic orbit \(\gamma\) of the shift \(\sigma_p\) on \(T_p\).
\item A choice of framing on each component of \(\beta\) (blackboard framing plus an explicit integer correction \(\mathrm{wr}(\beta) + f_k\) for the \(k\)-th strand).
\item An identification of the closure \(\overline{\beta}\) (in \(S^3\)) with a Seifert fibered 3-manifold \(M_\beta\) whose Seifert invariants \((\alpha_i, \beta_i)_{i=1}^r\) and Euler number \(e(M_\beta)\) are written down explicitly in terms of the combinatorial data of \(\gamma\).
\end{enumerate}
The Reshetikhin--Turaev invariant \(\mathrm{RT}_{\mathcal{C}_p}(\overline{\beta})\) is then evaluated by the general Seifert formula (Kirby--Melvin \cite{KirbyMelvin}, Turaev \cite{TuraevBook}):
\[
\mathrm{RT}_{\mathcal{C}_p}(M) = \frac{1}{D_p^{2g-2+r}} \sum_{j_1,\dots,j_r} \Bigl( \prod_{\ell=1}^r d_{j_\ell} \Bigr) S_{0j_1}\cdots S_{0j_r} \prod_{k=1}^s \theta_{j_k}^{m_k} \exp\bigl(2\pi i \, e \, h(\mathbf{j})\bigr),
\]
where \(h(\mathbf{j})\) is the explicit quadratic form in the Casimir values \(j_k(j_k+1)\) determined by the Seifert invariants, and the exponents \(m_k\) incorporate both the framing of \(\beta\) and the twisting parameters \((\alpha_i,\beta_i)\).

For the reduction to a pure power \(p^{-n}\) to hold after all normalizations, the state sum must collapse (via repeated Verlinde fusion rules and \(S\)-matrix orthogonality) to a \emph{different} matrix element than the one obtained in the canonical construction. Concretely, the current identification yields a reduction of the schematic form
\[
\mathrm{RT}_{\mathcal{C}_p}(\beta) = \frac{1}{D_p^{n-1}} \Biggl( \prod_{k=1}^n \theta_{j_k}^{e_k} \Biggr) G_p,
\]
with \(G_p = \sum_j d_j(p)\,\theta_j(p)\,S_{j0}\) (see Subsection~\ref{subsec:verlinde-reduction}). An alternative identification must produce instead a reduced expression whose leading term is a multiple of \(\tau_p\) (or an analogous weighted sum that equals \(\tau_p\) after the new framing phases are inserted). In other words, the Verlinde matrix element that survives after fusion must be altered from \(S_{j0}\) (or the column indexed by the trivial object) to a different column or a linear combination that compensates for the mismatch between \(G_p\) and \(\tau_p\).

\subsubsection{Explicit changes required in the braid encoding}
One concrete route is to replace the “natural” projection of \(\gamma\) onto the boundary \(\mathbb{P}^1(\mathbb{Q}_p)\) (which induces a single cyclic permutation on the \(n\) strands) by a different geometric realization. For instance:
\begin{itemize}
\item Choose a different base vertex \(v_0 \in V(T_p)\) and a different ordering of the \(p+1\) branches emanating from each vertex. The resulting permutation of strands may be a product of several disjoint cycles rather than a single \(n\)-cycle. The braid word \(\beta'\) then lies in a different conjugacy class in \(B_n\), and the closure \(\overline{\beta'}\) acquires different linking numbers between components.
\item Incorporate an auxiliary framing twist along each edge of the orbit \(\gamma\) (e.g., a full \(2\pi\) rotation of the tangent frame when traversing an edge of \(T_p\)). This adds an integer \(f_k\) to the writhe correction on the \(k\)-th strand, which propagates directly into the exponents \(m_k\) and \(e_k\) in the Verlinde formula.
\item Replace the Seifert fibration over a 4-punctured sphere by a fibration over a higher-genus orbifold or over a sphere with a different number \(r\) of exceptional fibers. The new Seifert invariants \((\alpha_i',\beta_i',e')\) must be chosen so that the quadratic form \(h(\mathbf{j})\) and the twisting exponents \(m_k\) exactly cancel the oscillatory factors in \(G_p\) while preserving the overall normalization \(D_p^{-(n-1)}\).
\end{itemize}
Each of these modifications changes the explicit Seifert data \((\alpha_i,\beta_i,e)\) and therefore the precise linear combination of \(S\)-matrix entries that survives after applying the Verlinde fusion rules \(N_{ab}^c = \sum_d S_{ad}S_{bd}S_{cd}/S_{0d}\) repeatedly. The surviving term is no longer forced to be the \(j'=0\) column; it may involve \(S_{j\ell}\) for a nontrivial \(\ell\) determined by the new exceptional-fiber multiplicities.

\subsubsection{Rigorous verification obligations}
Any such alternative identification must satisfy the following checklist (all steps must be carried out with explicit formulas, no deferred details):
\begin{enumerate}
\item Prove that the new braid word \(\beta'\) is independent of the choice of base vertex up to isotopy in \(S^3\) (Lemma analogous to Lemma~4.1 in the canonical case).
\item Compute the Seifert invariants \((\alpha_i',\beta_i',e')\) explicitly as functions of \(p\) and the combinatorial type of \(\gamma\) (e.g., via the continued-fraction expansion of the slope induced by the periodic point on \(\mathbb{P}^1(\mathbb{Q}_p)\)).
\item Substitute these invariants into the general Kirby--Melvin formula and reduce the state sum step-by-step, exhibiting every application of the Verlinde formula and every orthogonality relation \(\sum_j S_{0j}S_{j\ell} = D_p \delta_{0\ell}\).
\item Track every framing phase \(\theta_j^{e_k}\) (including the new writhe corrections) and show that they cancel against the exponential factor \(\exp(2\pi i e' h(\mathbf{j}))\) up to an overall root of unity of modulus \(1\).
\item Prove that the resulting reduced sum equals \(\tau_p\) (or a phase-adjusted version thereof) after division by \(D_p^{n-1}\), thereby obtaining \(\mathrm{RT}_{\mathcal{C}_p}(\beta') = p^{-n}\).
\end{enumerate}

\subsubsection{Why this route is currently open}
No explicit alternative braid word or set of Seifert invariants satisfying the above checklist has been constructed. The canonical encoding of Section~\ref{sec:braid} yields a single \(n\)-cycle whose closure is a Seifert manifold over the 4-punctured sphere with the invariants that produce precisely the Gauss sum \(G_p\). Any deviation requires a new geometric argument relating the \(p\)-adic shift \(\sigma_p\) to a different permutation representation on the strands, together with a new proof that the resulting 3-manifold is indeed Seifert fibered (which is not automatic for arbitrary braids). Numerical experiments with small \(p\) and \(n\) (Appendix~C) confirm that random modifications of the braid word do not restore the identity \(\mathrm{RT}_{\mathcal{C}_p}(\beta) = p^{-n}\).

\begin{note}[Conditional dependence on Section~\ref{subsec:verlinde-reduction}]
The discussion above is entirely conditional upon the gaps identified in Subsection~\ref{subsec:verlinde-reduction}. Until a concrete alternative identification is supplied together with the full checklist above, every claim in Sections~\ref{sec:proof}--\ref{sec:verification} and all later material of the paper (local Fredholm determinants, global adelic operator, functional equation, spectral reality, zero counting) remains conditional. No unconditional assertion is made in the present subsection.
\end{note}

This subsection records the precise mathematical requirements that any future correction of the braid encoding must meet. It may serve as a blueprint for a revised Phase~2 of the larger program while preserving full rigor at every step.